% Build the layer, Chapter 2, fourth build: a revocation that reaches every gate.
% Inserted after the third build (\S\ref{sec:anchor-build-taken-back}),
% whose gate it builds on. The design it builds part of is
% \S\ref{sec:revocation}'s: "Spreading revocations by gossip".
%
% No code or output in this file is typed by hand. \pdexcerpt, \pdrun and
% \pdrunlines print files under excerpts/, which the kernel's
% tools/excerpts.mjs cuts from its source and its recorded runs (see
% figures/pd-listing.tex). To change a listing, change the kernel and cut
% again. The prose names functions, never line numbers, so that a listing
% which moves in its file does not make a sentence wrong.
%
% \pdlistingroom{n} before a listing of n lines keeps it on one page. When a
% listing changes length, change the number with it.
\subsection{Build the layer: spreading a revocation by gossip}
\label{sec:anchor-build-told}

\noindent\emph{Express lane: a gate that has not been informed of a
revocation still accepts the card. Gates that gossip must not add to their lists
whatever they receive. The gate that works sends the chain the card was
revoked on, and each gate verifies it for itself. Among five gates a
revocation takes $2.37$ rounds on average to reach them all, and until it
reaches a gate the card is accepted there.}

\S\ref{sec:anchor-build-taken-back} built one gate, and said so. A harbor
has several, and \S\ref{sec:revocation} says how they should learn of a
revocation: each exchanges notices with another, chosen at random,
round after round. That section gives the time this takes as about
$\log D$ rounds for $D$ gates and marks the figure as unproved. This
section builds the gossip in \pdkernelname{}, and then counts the rounds.

\subsubsection{Break: a gate that was not informed}

Two gates trust one harbor. Call the one where a card is revoked the near
gate, and the other the far gate. Each is the gate of
Listing~\ref{alg:anchor-taken-gate}. Heron's key has leaked to Wren, as
before, and Alice revokes his card at the near gate.

\pdlistingroom{23}
\pdrunlines[caption={A card revoked at one gate is accepted at the
other.},label={alg:anchor-told-break}]{two-gates-local}{1,9-30}{1 and 9\textendash 30}

The near gate refuses the card and the far gate accepts it. A minute later
the gates exchange notices, and nothing changes, because these gates do not
gossip: neither has a way to send what is on its list. Wren only has to go
to the other gate.

\subsubsection{A first repair: send the list}

A gate's revocation list is a list of entries, and
\S\ref{sec:anchor-build-taken-back} said what an entry is: a card's id,
and the card's expiry. So let each gate send its entries, and
add the ones it receives.

In Listing~\ref{alg:anchor-told-hearsay}, \texttt{notices} is what a gate
sends, and \texttt{receive} is what it does with another gate's notices. Run
the break again with these gates and it is fixed: once the two have exchanged
notices, the far gate refuses Heron's card as the near gate does.

\pdlistingroom{16}
\pdexcerpt[caption={The first attempt. A gate adds whatever it
receives.},label={alg:anchor-told-hearsay}]{ch02d-unverified}

\subsubsection{Break again: a notice from someone who is not a gate}

\texttt{receive} accepts whatever arrives. Nothing in an entry shows that a
gate sent it, or that anyone ever held the card it names. Wren can send
entries too. She knows the id of Heron's card, because she has the card.
She knows the id of Finch's as well, without ever seeing it: the harbor
numbers its own cards, and his is \texttt{card-2}.

She sends the far gate an entry for each, and each does its own harm
(Listing~\ref{alg:anchor-told-rumour}). The entry for Finch's card has
him refused, on a card nobody revoked. The entry for Heron's card says
that the card expires at $t=180000$, which is now. When the true entry
comes from the near gate, the far gate already has one for that id and
keeps the one it has. A millisecond later the rule of
Listing~\ref{alg:anchor-taken-spent} allows that entry to be pruned, since
the card it describes has expired, and Heron's card is accepted again at
the far gate with seven minutes left.

\pdlistingroom{24}
\pdrunlines[caption={A gate that adds what it receives, and two entries
sent by someone who is not a gate.},label={alg:anchor-told-rumour}]{false-notice-unverified}{1,13-35}{1 and 13\textendash 35}

This is the early pruning of \S\ref{sec:anchor-build-taken-back}, and this
time someone outside the harbor caused it. One line of the listing is easy
to miss: when the gates exchange notices, the near gate adds one card to its
list. That is Finch's. The far gate has forwarded the false entry as its
own.

The listing also shows something Wren tried that did not work, and it
leads to the fix. She sent Heron's real cards as a notice, with the expiry
of his card changed. This gate ignored them, because it expects entries.

\subsubsection{Fix: send the chain}

An entry is a claim about a card. The chain is the card itself, with its
ancestors. So the notice becomes the chain a card was revoked on,
and a gate that receives a chain treats it as it would a chain presented
to it directly.

\pdlistingroom{27}
\pdexcerpt[caption={The gate that gossips. The notice is the chain, and
each gate verifies it.},label={alg:anchor-told-tell}]{ch02d-verified}

\texttt{revoke} is the third build's, and it also keeps the chain.
\texttt{receive} passes each chain it has not seen to \texttt{revoke}. That
is all of it. The core's check of a revoked chain
(\S\ref{sec:anchor-build-taken-back}) now runs at every gate that receives
the notice, and each gate adds its own entry, with the expiry that the core
read from the card. No gate trusts what another gate says, so no gate
needs to know who is sending.

Two earlier results make this safe to do. Since
\S\ref{sec:anchor-build-holder} a card is not a secret: without its
holder's key it grants nothing, so sending cards from gate to gate gives
nobody access. And a card's signature covers its claims. Wren's changed
card has the real card's id, because the id is the hash of the signature
and she left the signature alone. The signature no longer matches the
claims, and the core refuses the chain.

% A page may break here. See the note in anchor-build-the-verifier.tex.
\pagebreak[0]
\keyidea[Send the evidence, not the verdict]{A gate that is told ``this
card was revoked'' has to decide whom to trust. A gate that is sent
the chain has nothing to decide. It runs the check it would have run
anyway.}

\subsubsection{Verify: the same runs}

The gate that was not informed, first
(Listing~\ref{alg:anchor-told-heard}).

\pdlistingroom{23}
\pdrunlines[caption={Two gates that gossip. The far gate refuses the card
once it has received the notice.},label={alg:anchor-told-heard}]{two-gates-core}{1,9-30}{1 and 9\textendash 30}

Look at Wren's first request at the far gate. Before the gates
exchange notices, the far gate admits her, under this design as
under every other. A gate cannot refuse for a reason that has not reached
it. The verdict says so, and the rest of this section measures how long
that lasts.

Then the false notice (Listing~\ref{alg:anchor-told-refused}). Wren's
entries are not chains, and her changed chain fails the check. Finch is
admitted. The far gate receives the true chain from the near gate, and its
entry stays until the card expires.

\pdlistingroom{20}
\pdrunlines[caption={The false notice, against the fix. Nothing Wren
sends is added.},label={alg:anchor-told-refused}]{false-notice-core}{1,15-33}{1 and 15\textendash 33}

This build added no check of its own, so it has no guard of its own to
remove. The guard it relies on is the signature check of
\S\ref{sec:anchor-build}. Build the core without that
(Listing~\ref{alg:anchor-told-mutant}).

\pdlistingroom{15}
\pdrunlines[caption={The fix with the signature check removed. The
changed chain is added, and the revoked card is accepted
again.},label={alg:anchor-told-mutant}]{false-notice-unsigned}{1,19-32}{1 and 19\textendash 32}

\subsubsection{How long a revocation takes to spread}

What is left is the line in Listing~\ref{alg:anchor-told-heard} where the
far gate admits Wren. How long does a gate stay uninformed?
\S\ref{sec:revocation} describes the rounds, and here they are as code.

\pdlistingroom{17}
\pdexcerpt[caption={One round. Every gate picks another, and the two
exchange notices.},label={alg:anchor-told-round}]{ch02d-round}

The lists are read once, before any gate receives anything. So a notice
travels one gate further in a round and no further: a gate cannot forward
in round three what it received in round three.

Three gates can be done by hand. One is informed. In the first round it
picks one of the other two and informs it. The remaining gate is also
informed in that round if the gate it picks is the one that began
informed, and it picks that gate half the time. If it does not, it picks
an informed gate in the next round, whichever it picks. So the notice takes
one round or two, each half the time, and $\tfrac32$ rounds on average.
Exercise~\ref{ex:anchor-told-three} asks for the wait of a single gate.

For more gates the same reasoning needs one probability, and the kernel
computes it exactly, in fractions.

\pdlistingroom{25}
\pdexcerpt[caption={The probability that a named set of gates all end a
round still uninformed.},label={alg:anchor-told-quiet}]{ch02d-uninformed}

Only the number of informed gates matters, because the gates are
alike, and that number never falls. From \texttt{stayUninformed} the kernel
gets the probability of every increase, and from those the expected number
of rounds (Listing~\ref{alg:anchor-told-spread}). For three to six gates it
also tries every choice of peers the gates could make, and the two agree.

The column \texttt{all} is the wait until every gate is informed, and
\texttt{one} is the wait of a gate named in advance, which is how long a
revoked card is still accepted at that gate. With five gates they are $2.37$
and $1.68$ rounds.

\pdlistingroom{18}
\pdrunlines[caption={The expected rounds, computed
exactly.},label={alg:anchor-told-spread}]{spread}{1,11-27}{1 and 11\textendash 27}

From 4 gates to 64, each doubling adds about one round,
and a little less each time. That is the $\log D$ of
\S\ref{sec:revocation}, with its constants. The last remark in the
listing compares the exact figures with the growth that is proved for large
networks, $\log_3 D + \log_2 \ln D$ up to a constant~\pdcite{doerr2017rumor}:
from 8 gates to 64 the two differ by a steady sixth of a round.

A model can be wrong about the code it describes. So the kernel also runs
its gates, the ones of Listing~\ref{alg:anchor-told-tell}, in the rounds
of Listing~\ref{alg:anchor-told-round}, 400 times for each of six sizes,
and counts the rounds. Every run starts from a seed, so the counts are the
same each time they are made.

The average of 400 runs is not the true figure: another 400 runs would
give another average. Runs are independent, so their mean varies around
the true figure by about one standard error, which is the standard
deviation of the runs divided by the square root of their number, and it
lies within $1.96$ of those 19 times in 20.
Listing~\ref{alg:anchor-told-gossip}
gives that interval for each size as \texttt{low} and \texttt{high}. If the
model describes the gates, its figure, \texttt{exact}, should fall inside
the interval in most rows.

\pdlistingroom{12}
\pdrunlines[caption={400 runs of the gates at each of six sizes, beside
the model.},label={alg:anchor-told-gossip}]{gossip}{1,12-22}{1 and 12\textendash 22}

It does in all six. Figure~\ref{fig:anchor-gossip-rounds} draws them,
with the same runs made again while one exchange in four is lost. Loss
slows the spread and does not stop it.

% Brief (grammar: measured against exact, two panels; atlas
%   V/fig:anchor-gossip-rounds).
% Reader question: how many rounds does news of a card taken back need to
% reach every gate, and can the count be trusted?
% Claim: the rounds grow by about one for each doubling of the gates; runs
% of the kernel's gates agree with the exact figures; and the one fit that
% looks loose is what chance does about once in 72 tries.
% Evidence: exact expected rounds for 2 to 64 gates; 400 seeded runs of the
% kernel's gates at each of six sizes, with no exchange lost and with one
% in four lost; the fit of the five-gate rounds made for 200 other seeds
% [internal, harbor spread, harbor gossip, harbor fits; sessions
% spread.txt, gossip.txt, gossip-loss.txt and fits.txt, plotted as printed].
% Must distinguish: worked out from counted; every gate from a given gate;
% no loss from loss; the 200 seeds from this Book's own.
% Grammar chosen: on the left, gates on a doubling scale, so that a steady
% gain for each doubling is a straight line; a line for what is worked out,
% dots with interval bars for what was counted. On the right, ten bars
% against the level chance alone gives.
% Grammar rejected: a linear scale of gates, which crowds 2 to 8 into a
% corner; a single bar for "the fit passed", which hides how often chance
% alone gives a fit this loose.
% Five-second test: do the dots sit on the lines?
\begin{figure}[H]
\centering
\pdfigurehue{pdcobalt}%
\begin{tikzpicture}[pd figure]
  \pgfplotsset{
    counted/.style={only marks,mark size=1.3pt,
      error bars/y dir=both,error bars/y explicit,
      error bars/error bar style={line width=.5pt},
      error bars/error mark options={rotate=90,mark size=1.1pt,line width=.5pt}},
  }
  % The tables are the recorded runs, set as they were printed: the first
  % line is the command, and every remark starts with '#'.
  \pgfplotstableset{
    as printed/.style={skip first n=1,x=gates},
    bars/.style={y=mean,y error plus expr=\thisrow{high}-\thisrow{mean},
      y error minus expr=\thisrow{mean}-\thisrow{low}},
  }
  \begin{axis}[pd axis,at={(0,0)},width=4.2cm,height=4.9cm,
    xmode=log,log basis x=2,xmin=1.7,xmax=76,ymin=0,ymax=8.2,
    xtick={2,4,8,16,32,64},xticklabels={2,4,8,16,32,64},
    ytick={0,2,4,6,8},
    title={rounds until the gates have heard},xlabel={gates},ylabel={rounds}]
    \addplot[pd guide,mark=none] table[as printed,y=one] {excerpts/spread.txt};
    \addplot[pd muted series,mark=none] table[as printed,y=exact] {excerpts/gossip-loss.txt};
    \addplot[counted,draw=pdinkmuted!70,mark options={fill=pdinkmuted!70,draw=pdpage,line width=.3pt}]
      table[as printed,bars] {excerpts/gossip-loss.txt};
    \addplot[draw=pdfocus,line width=.9pt,mark=none] table[as printed,y=all] {excerpts/spread.txt};
    \addplot[counted,draw=pdink,mark options={fill=pdink,draw=pdpage,line width=.3pt}]
      table[as printed,bars] {excerpts/gossip.txt};
    \draw[draw=pdfocus,line width=.3pt] (axis cs:64,6.12) -- (axis cs:64,7.25);
    \node[pd focus label,anchor=south east] at (axis cs:76,7.2) {every gate};
    \node[pd label,text=pdinkmuted,anchor=south west] at (axis cs:2.2,5.6) {one in four lost};
    \node[pd label,anchor=north east] at (axis cs:76,2.45) {a given gate};
  \end{axis}
  \begin{axis}[pd axis,at={(5.45cm,0)},width=3.4cm,height=4.9cm,
    xmin=0,xmax=1,ymin=0,ymax=32,
    xtick={0,0.5,1},xticklabels={0,.5,1},ytick={0,10,20,30},
    title={the fit, for 200 other seeds},xlabel={chance of a fit that loose},ylabel={seeds}]
    \addplot[ybar,bar width=.1,bar shift=0pt,draw=pdpage,line width=.3pt,fill=pdinkmuted!55]
      table[skip first n=1,x expr=\thisrow{chance}+0.05,y=seeds] {excerpts/fits.txt};
    \addplot[pd guide,mark=none] coordinates {(0,20) (1,20)};
    % This Book's own seed: its chance is the first number of the run's
    % closing remark, held there by the numbers test.
    \draw[pd focus rule] (axis cs:0.0138,0) -- (axis cs:0.0138,29.5);
    \node[pd focus label,anchor=west] at (axis cs:0.04,29.3) {this Book's seed};
  \end{axis}
\end{tikzpicture}
\caption{\textbf{News reaches every gate in about a round for each doubling
of the gates, and runs of the gates agree with the arithmetic.} On the left,
lines are exact and dots are the means of 400 runs with 95\% intervals. The
lowest line is the wait of one gate named beforehand. On the right, the
chance of the five-gate fit for seeds 1 to 200, by tenths; chance alone puts
20 in each, the dotted level. [internal, \texttt{harbor spread}, \texttt{harbor gossip},
\texttt{harbor fits}]}
\label{fig:anchor-gossip-rounds}
\end{figure}


A partition does stop it. Divide the gates into two halves that cannot
reach each other for five rounds (Listing~\ref{alg:anchor-told-cut}).

\pdlistingroom{4}
\pdrunlines[caption={The gates partitioned in two for five
rounds.},label={alg:anchor-told-cut}]{gossip-partition}{1,14-16}{1 and 14\textendash 16}

\texttt{least} is the fewest rounds any of 400 runs took, and it is 6:
no run finished while the partition lasted. \S\ref{sec:revocation} calls
its figure an average and not a deadline, and this is the reason.

A mean can agree while the distribution behind it is wrong. So take five
gates, make 4{,}000 runs, and count how many took each number of rounds
(Listing~\ref{alg:anchor-told-each}).

\pdlistingroom{12}
\pdrunlines[caption={The rounds five gates took in 4{,}000 runs, against
the runs the model expects.},label={alg:anchor-told-each}]{gossip-each}{1,11-21}{1 and 11\textendash 21}

The model says one run in 64 ends in the first round. The informed gate
picks one other gate and informs it, and each of the remaining three must
pick the informed gate, with a probability of a quarter each. $4000/64$ is
the $62.50$ under \texttt{expected}. The runs gave 45. Chi-square adds up,
over the rows, the squared difference between \texttt{observed} and
\texttt{expected}, divided by \texttt{expected}, and the sum is $10.65$. A
sum that size is unlikely: when runs do follow the model, the probability
of one that large or larger is $0.0138$, about once in 72 tries.

That number was not chosen. The seed is the one every count in this book
uses. The significance level, $0.01$, was set before any run. The size was
fixed at 4{,}000 after a first look at half as many, which had also come out
unusual, and the larger run was made once. $0.0138$ is above the level, and
close to it. A figure like this is a reason to measure again with data
independent of the first measurement, and it is no reason to change the
seed. So the kernel
makes the same fit for seeds 1 to 200. The right of
Figure~\ref{fig:anchor-gossip-rounds} shows where their p-values fall. If
the model were wrong they would cluster near zero. They are close to
uniform, and 3 of the 200 are smaller than this book's, where a uniform
distribution expects $2.8$.
The fit of Listing~\ref{alg:anchor-told-each} is what one draw in 72
looks like.
\pdscope{The exact figures assume that every gate can reach every other
and picks uniformly among them, in rounds that begin together. The runs
cover 3 to 32 gates in one process on one clock. Nothing here is a
machine-checked proof, and nothing bounds the time under a partition.}

\begin{pdboundary}
\setlength{\parskip}{4pt plus 1pt}
A gate that is not yet informed accepts the card. Gossip shortens that
time, to $1.68$ rounds on average for one gate among five, and no design
of this kind removes it. A card whose misuse cannot wait that long needs a
short lifetime, which is \S\ref{sec:anchor-build}'s tool and not this one's.

Sending chains has a cost that \S\ref{sec:anchor-build-taken-back}
warned of. The gate still does not check who is revoking a card, and now
every gate that receives a chain holds the ancestors of the revoked card,
as does anyone who can read the network traffic. Each of those cards is a
chain of its own, and whoever holds it can have it revoked at every gate
(Exercise~\ref{ex:anchor-told-above}). Nobody gains a right this way, and
a holder can be locked out until a new card is signed. Deciding who may
revoke a card is the next thing this kernel needs, and
\S\ref{sec:anchor-build-warranted} builds it.

A chain also reveals more than an id does: who holds which rights,
delegated from whom. \S\ref{sec:revocation} suggests sending salted hashes
of ids when that matters, and a hash cannot be verified the way a chain
can. The rounds here begin together on one clock, and every gate can reach
every other. Real gates differ in both ways, and the figures above are for
the simplest case.
\end{pdboundary}
